% GENERATED FILE. Edit canonical lessons, then run npm run generate:books.
% canonical-source-hash: fnv1a64:873b0b4c9c43026f
% canonical-lessons: UR-C182-thakna, UR-C182-aram-karna, UR-C182-letna, UR-C182-khara-hona, UR-C182-intizar-karna

\chapter{To get tired, To rest, To lie down, To stand up, To wait}
\label{ch:ur-to-get-tired-to-rest-to-lie-down-to-stand-up-to-wait}
\hlchaptermodality{182}

\begin{chapteropening}
\textbf{By the end of this chapter:} \emph{I can say to get tired, to rest, to lie down, to stand up and to wait.}

Complete the last lesson of chapter 182: I can say to get tired, to rest, to lie down, to stand up and to wait.
\end{chapteropening}

\section[\texorpdfstring{\ur{تھکنا} (thaknā)}{تھکنا (thaknā)}]{\ur{تھکنا} (thaknā) — to get tired}

\label{lesson:UR-C182-thakna}



\begin{quote}
Before the new one: say the Urdu for to weigh, then the Urdu for to check.
\end{quote}

\subsection*{You'll want to know: \ur{تھکنا}}
\textbf{\ur{تھکنا}} — \emph{thaknā} — \textquotedblleft{}to get tired\textquotedblright{}.

\begin{cousinweb}
\textbf{\ur{تھکنا}} (\emph{thaknā}) is a tadbhava, a word inherited and reshaped by sound change. It comes from \emph{*sthakk}, \textquotedblleft{}stop, halt\textquotedblright{}, a form CDIAL reconstructs from its descendants (CDIAL 13737). CDIAL calls a link with \emph{sthā-}, \textquotedblleft{}to stand\textquotedblright{}, probable but not clear.
\end{cousinweb}

\begin{grammarlens}[title={What you've built}]
One more word for this chapter.
\end{grammarlens}

\subsection*{Your turn}
\begin{itemize}
\raggedright
  \item \emph{Say it:} \emph{thaknā}
  \item \emph{Say it:} \emph{thaknā}, once more
  \item \emph{Say it:} say \emph{thaknā}
  \item \emph{Recall:} say the Urdu for to weigh, then the Urdu for to check, then say \emph{thaknā} again
\end{itemize}

\begin{tcolorbox}[breakable,colback=teal!4,colframe=teal!35!black,title={Before you move on}]
What does \ur{تھکنا} mean? (\textquotedblleft{}To get tired\textquotedblright{}.) Say it once more.
\end{tcolorbox}
\section[\texorpdfstring{\ur{آرام} \ur{کرنا} (ārām karnā)}{آرام کرنا (ārām karnā)}]{\ur{آرام} \ur{کرنا} (ārām karnā) — to rest}

\label{lesson:UR-C182-aram-karna}



\begin{quote}
Before the new one: say the Urdu for to check, then the Urdu for to get tired.
\end{quote}

\subsection*{You'll want to know: \ur{آرام} \ur{کرنا}}
\textbf{\ur{آرام} \ur{کرنا}} — \emph{ārām karnā} — \textquotedblleft{}to rest\textquotedblright{}.

\begin{cousinweb}
\textbf{\ur{آرام} \ur{کرنا}} (\emph{ārām karnā}) joins \emph{ārām}, \textquotedblleft{}rest\textquotedblright{}, usually traced to Persian \emph{ārām}, and \emph{karnā}, \textquotedblleft{}to do\textquotedblright{}, from Sanskrit \emph{karoti}, \textquotedblleft{}does\textquotedblright{} (CDIAL 2814).
\end{cousinweb}

\begin{grammarlens}[title={What you've built}]
One more word for this chapter.
\end{grammarlens}

\subsection*{Your turn}
\begin{itemize}
\raggedright
  \item \emph{Say it:} \emph{ārām karnā}
  \item \emph{Say it:} \emph{ārām karnā}, once more
  \item \emph{Say it:} say \emph{ārām karnā}
  \item \emph{Recall:} say the Urdu for to check, then the Urdu for to get tired, then say \emph{ārām karnā} again
\end{itemize}

\begin{tcolorbox}[breakable,colback=teal!4,colframe=teal!35!black,title={Before you move on}]
What does \ur{آرام} \ur{کرنا} mean? (\textquotedblleft{}To rest\textquotedblright{}.) Say it once more.
\end{tcolorbox}
\section[\texorpdfstring{\ur{لیٹنا} (leṭnā)}{لیٹنا (leṭnā)}]{\ur{لیٹنا} (leṭnā) — to lie down}

\label{lesson:UR-C182-letna}



\begin{quote}
Before the new one: say the Urdu for to get tired, then the Urdu for to rest.
\end{quote}

\subsection*{You'll want to know: \ur{لیٹنا}}
\textbf{\ur{لیٹنا}} — \emph{leṭnā} — \textquotedblleft{}to lie down\textquotedblright{}.

\begin{cousinweb}
\textbf{\ur{لیٹنا}} (\emph{leṭnā}) is a tadbhava, a word inherited from Sanskrit and reshaped by sound change. It comes from \emph{leṭyati}, \textquotedblleft{}sleeps\textquotedblright{} (CDIAL 11109). The Sanskrit verb is known only from a grammarians' list, and CDIAL thinks it probably of non-Aryan origin.
\end{cousinweb}

\begin{grammarlens}[title={What you've built}]
One more word for this chapter.
\end{grammarlens}

\subsection*{Your turn}
\begin{itemize}
\raggedright
  \item \emph{Say it:} \emph{leṭnā}
  \item \emph{Say it:} \emph{leṭnā}, once more
  \item \emph{Say it:} say \emph{leṭnā}
  \item \emph{Recall:} say the Urdu for to get tired, then the Urdu for to rest, then say \emph{leṭnā} again
\end{itemize}

\begin{tcolorbox}[breakable,colback=teal!4,colframe=teal!35!black,title={Before you move on}]
What does \ur{لیٹنا} mean? (\textquotedblleft{}To lie down\textquotedblright{}.) Say it once more.
\end{tcolorbox}
\section[\texorpdfstring{\ur{کھڑا} \ur{ہونا} (khaṛā honā)}{کھڑا ہونا (khaṛā honā)}]{\ur{کھڑا} \ur{ہونا} (khaṛā honā) — to stand up}

\label{lesson:UR-C182-khara-hona}



\begin{quote}
Before the new one: say the Urdu for to rest, then the Urdu for to lie down.
\end{quote}

\subsection*{You'll want to know: \ur{کھڑا} \ur{ہونا}}
\textbf{\ur{کھڑا} \ur{ہونا}} — \emph{khaṛā honā} — \textquotedblleft{}to stand up\textquotedblright{}.

\begin{cousinweb}
\textbf{\ur{کھڑا} \ur{ہونا}} (\emph{khaṛā honā}) is \textquotedblleft{}to become standing\textquotedblright{}. CDIAL lists \emph{khaṛā}, \textquotedblleft{}standing erect\textquotedblright{}, under Sanskrit \emph{khaḍaka}, \textquotedblleft{}a bolt, a post\textquotedblright{}, for which it supposes an older sense \textquotedblleft{}erect\textquotedblright{} (CDIAL 3784). \emph{honā}, \textquotedblleft{}to be\textquotedblright{}, comes from Sanskrit \emph{bhavati}, \textquotedblleft{}becomes, is\textquotedblright{} (CDIAL 9416).
\end{cousinweb}

\begin{grammarlens}[title={What you've built}]
One more word for this chapter.
\end{grammarlens}

\subsection*{Your turn}
\begin{itemize}
\raggedright
  \item \emph{Say it:} \emph{khaṛā honā}
  \item \emph{Say it:} \emph{khaṛā honā}, once more
  \item \emph{Say it:} say \emph{khaṛā honā}
  \item \emph{Recall:} say the Urdu for to rest, then the Urdu for to lie down, then say \emph{khaṛā honā} again
\end{itemize}

\begin{tcolorbox}[breakable,colback=teal!4,colframe=teal!35!black,title={Before you move on}]
What does \ur{کھڑا} \ur{ہونا} mean? (\textquotedblleft{}To stand up\textquotedblright{}.) Say it once more.
\end{tcolorbox}
\section[\texorpdfstring{\ur{انتظار} \ur{کرنا} (intizār karnā)}{انتظار کرنا (intizār karnā)}]{\ur{انتظار} \ur{کرنا} (intizār karnā) — to wait}

\label{lesson:UR-C182-intizar-karna}



\begin{quote}
Before the new one: say the Urdu for to lie down, then the Urdu for to stand up.
\end{quote}

\subsection*{You'll want to know: \ur{انتظار} \ur{کرنا}}
\textbf{\ur{انتظار} \ur{کرنا}} — \emph{intizār karnā} — \textquotedblleft{}to wait\textquotedblright{}.

\begin{cousinweb}
\textbf{\ur{انتظار} \ur{کرنا}} (\emph{intizār karnā}) joins \emph{intizār}, \textquotedblleft{}waiting\textquotedblright{}, and \emph{karnā}, \textquotedblleft{}to do\textquotedblright{}. \emph{intizār} is Arabic \emph{intiẓār}, \textquotedblleft{}waiting, looking out for\textquotedblright{}, which reached Urdu through Persian (Lane, root \emph{n-ẓ-r}); \emph{karnā} comes from Sanskrit \emph{karoti}, \textquotedblleft{}does\textquotedblright{} (CDIAL 2814).
\end{cousinweb}

\begin{grammarlens}[title={What you've built}]
One more word for this chapter.
\end{grammarlens}

\subsection*{Your turn}
\begin{itemize}
\raggedright
  \item \emph{Say it:} \emph{intizār karnā}
  \item \emph{Say it:} \emph{intizār karnā}, once more
  \item \emph{Say it:} say \emph{intizār karnā}
  \item \emph{Recall:} say the Urdu for to lie down, then the Urdu for to stand up, then say \emph{intizār karnā} again
\end{itemize}

\begin{tcolorbox}[breakable,colback=teal!4,colframe=teal!35!black,title={Before you move on}]
What does \ur{انتظار} \ur{کرنا} mean? (\textquotedblleft{}To wait\textquotedblright{}.) Say it once more.
\end{tcolorbox}
